\documentclass[12pt,a4paper,oneside]{article}
\usepackage[brazil]{babel}
\usepackage[utf8]{inputenc}
\usepackage[dvipsnames]{xcolor}
\usepackage{listings}
\usepackage{olymp}

\setcounter{problem}{0}

\begin{document}

\begin{problem}{Horário}{horario}{2}
 
Faça um programa para ler um horário no formato 24 horas e escrevê-lo no formato HH:MM AM/PM. Para escrever o horário você deve utilizar obrigatoriamente uma função com o seguinte protótipo:

\lstset{language=C++,
                basicstyle=\ttfamily,
                keywordstyle=\color{Mulberry}\ttfamily,
                stringstyle=\color{Maroon}\ttfamily,
                commentstyle=\color{YellowGreen}\ttfamily,
                morecomment=[l][\color{magenta}]{\#},
                basewidth = {.48em}
}
\begin{lstlisting}
void escreve (Horario h);
\end{lstlisting}

onde \texttt{Horario} é um tipo definido da seguinte forma:

\lstset{language=C++,
                basicstyle=\ttfamily,
                keywordstyle=\color{Mulberry}\ttfamily,
                stringstyle=\color{Maroon}\ttfamily,
                commentstyle=\color{YellowGreen}\ttfamily,
                morecomment=[l][\color{magenta}]{\#},
                basewidth = {.48em}
}
\begin{lstlisting}
struct Horario {
  int horas, minutos;
};
\end{lstlisting}


%\Notes
%
%Observações, se tiver.

\InputFile

A primeira linha da entrada contém um número $N$, indicando que haverá $N$ casos de teste. As $N$ linhas seguintes contêm um caso de teste cada, que são dois inteiros $H$ e $M$, respectivamente o número de horas e minutos. Restrições: $0 \leq H \leq 23, 0 \leq M \leq 59$.

\OutputFile

Seu programa deve gerar $N$ linhas de saída, uma para cada caso de teste da entrada, com o horário no formato desejado.

% \Notes
% Preste atenção aos tipos das variáveis, pois $F(90) \approx 2^{61}$.

\Example

\begin{example}
\exmp{
6
10 20
22 20
12 0
0 0
12 5
13 0
}{
10:20 AM
10:20 PM
12:00 PM
12:00 AM
12:05 PM
01:00 PM
}%
\end{example}



%Comentários sobre o exemplo, se necessário.

\end{problem}

\end{document}
